% ============================================================
% f2014_hunan_arts_q18_dihedral.tex
% 2014 湖南卷（文科）第 18 题：二面角 α-MN-β 与菱形 ABCD
% 来源 PDF：普通高考/2014/2014湖南文.pdf
% 原页：第 1 页（A3 横排）第 3 栏下半，印面 887，第 18 题题干下方
% 结构：半平面 α 在棱 MN 下方，半平面 β 在棱 MN 上方，二者成 60°；
%       菱形 ABCD 在 β 内，A、B 在 MN 上，E 为 AB 中点；
%       DO⊥α 于 O。实线：菱形四边、DE、DO；虚线：MN 上被平面 DOE 遮住的一段。
% ============================================================
\begin{tikzpicture}[>=Stealth, line width=0.7pt, font=\small,
  x=1cm, y=1cm, line cap=round, line join=round]
  % 半平面 β
  \draw (-3.60,0) -- (-4.05,2.90) -- (3.15,2.90) -- (3.60,0);
  % 半平面 α
  \draw (-3.60,0) -- (-5.20,-2.40) -- (2.00,-2.40) -- (3.60,0);
  % 棱 MN：M→A→(D 的垂足) 实线，(垂足)→E 虚线，E→B→N 实线
  \draw (-3.60,0) -- (-1.50,0);
  \draw[dashed] (-1.50,0) -- (-0.30,0);
  \draw (-0.30,0) -- (3.60,0);
  % 菱形 ABCD 的四边与对角线 DE
  \draw (-2.30,0) -- (-1.50,2.50);
  \draw (-1.50,2.50) -- (2.50,2.50);
  \draw (2.50,2.50) -- (1.70,0);
  \draw (-0.30,0) -- (-1.50,2.50);
  % DO⊥α
  \draw (-1.50,2.50) -- (-1.50,-1.05);
  % 顶点与标注
  \node[left]       at (-3.75,0.05) {$M$};
  \node[right]      at (3.75,0.05) {$N$};
  \node[below]      at (-2.30,-0.22) {$A$};
  \node[below right] at (-0.22,-0.10) {$E$};
  \node[below]      at (1.70,-0.22) {$B$};
  \node[above left] at (-1.50,2.50) {$D$};
  \node[above right] at (2.42,2.50) {$C$};
  \node[below]      at (-1.50,-1.35) {$O$};
  \node at (-3.05,2.30) {$\beta$};
  \node at (-3.95,-1.55) {$\alpha$};
\end{tikzpicture}
